\begin{theorem}[Injective parent uniqueness from vertex and edge coverage]%
    \label{thm:peps_injective_covering_parent_ground_space}
    \lean{TNLean.PEPS.globalVertexTensorMap_virtualBondProduct,
        TNLean.PEPS.regionParentGroundSpace_eq_span_of_isVertexInjective_of_cover,
        TNLean.PEPS.ker_regionParentHamiltonian_eq_span_of_isVertexInjective_of_cover,
        TNLean.PEPS.finrank_regionParentGroundSpace_eq_one_of_isVertexInjective_of_cover}
    \leanok
    \uses{thm:peps_covering_parent_virtual_image,
        thm:peps_regional_virtual_conditions_bell,
        thm:peps_bell_product_physical_image}
    Let the site tensors on a finite graph be injective. Suppose a
    collection of regions contains every vertex in some region and
    contains the two endpoints of every edge together in some region.
    Then its common parent space is
    \[
        \mathcal G=\operatorname{span}\{\Psi(A)\}.
    \]
    Every finite positive interaction family with the genuine regional
    kernels has this same zero-energy space. If all bond dimensions
    are positive, the closed vector is nonzero and the ground space
    is one-dimensional. The span formula also holds when a bond
    dimension vanishes, in which case the span can be zero.

    This is the covering-region form of the virtual-pair argument in
    \cite[Section IV.C.1]{Cirac2021Matrix}. Vertex coverage is explicit,
    including for isolated vertices. The physical reconstruction
    correction is recorded in
    \cite{gap:cpgsv21_injective_parent_reconstruction}.
\end{theorem}

\begin{proof}\leanok
    Vertex coverage supplies the genuine site-image conditions, so
    every physical parent vector is the product-site image of a
    vector satisfying all actual virtual regional conditions.
    Edge coverage imposes the Bell condition on each separate virtual
    bond. Their common space is precisely the line spanned by the
    product Bell vector. Its image is the actual closed PEPS vector.
    Conversely the closed vector satisfies every regional condition.
    Positivity gives the Hamiltonian-kernel statement, and positive
    bond dimensions give nonvanishing and dimension one.
\end{proof}
